\section{Minimum gauge defects generate the full flat connection space}
\label{sec:minimum-gauge-defect-generation}
The filled port complex has $45$ vertices, $720$ edges, and $240$ edge-disjoint filled triangles, with coefficient module $A=C_3^5\cong\mathbb F_3^5$.  On one oriented triangle the scalar flat plane is
\[
  \{(a,b,c)\in\mathbb F_3^3:a+b+c=0\},
\]
which has rank two and is spanned by its six nonzero weight-two vectors.  Since the filled triangles partition all graph edges, these local planes sum directly.  Consequently
\[
  \dim Z^1(X;\mathbb F_3)=240\cdot2=480,
  \qquad
  \boxed{\dim Z^1(X;\mathbb F_3^5)=2400}.
\]
The block graph is connected, so
\[
  \dim B^1(X;\mathbb F_3^5)=(45-1)\cdot5=220,
\]
and therefore
\[
  \boxed{\dim H^1(X;\mathbb F_3^5)=2400-220=2180}.
\]
Thus the weight-two minimum defects do not merely witness the minimum nontrivial class: they span every flat $C_3^5$ connection, and every switching class has a representative assembled from minimum defects.  This is a generation theorem, not a unique-decoding or multi-fault-correction claim.
